\section{A rank criterion on the conductor grid}\label{v2:sec:rank}
The component classification has an exact rational-rank formulation. It is useful both as an alternative decision certificate and for clarifying why a power of two appears in the mate count.

Fix a nonempty subgroup union $A$, let $V_A$ be the nonzero-conductor vertex set, and write $\delta_{\s}$ for the coordinate row vector indexed by $\s\in V_A$. Define $R_A$ over $\mathbb Q$ as follows: its columns are indexed by $V_A$; for each edge $\{\s,\mathbf t\}$ of $\Gamma_A$ include the row $\delta_{\s}-\delta_{\mathbf t}$; and for each root or pure-prime marked vertex $\s$ include the row $\delta_{\s}$. Any orientation of the edge rows gives the same rank. Repeated anchoring rows may be omitted.

\begin{theorem}[Conductor-constraint rank criterion]\label{v2:thm:rank}
For every positive integer $N$ and every nonempty subgroup union $A\subseteq\mathbb Z_N$,
\[
 \dim_{\mathbb Q}\ker R_A=u_{\rm div}(A).
\]
Moreover, $\ker R_A\cap\{0,1\}^{V_A}$ is in bijection with the subgroup-union homometric mates of $A$. For $z$ in this intersection, define
\[
 G(\s)=\begin{cases}
       (1-2z_{\s})F(\s),&\s\in V_A,\\
       0,&\s\notin V_A.
       \end{cases}
\]
Equations~\eqref{ii:eq:inversealpha} and $b(\vv)=\sum_{\rvec\leq\vv}\alpha_{\rvec}$ construct the mate. Consequently
\[
 \#\{B:\ B\text{ is a subgroup union and }C_B=C_A\}
   =2^{|V_A|-\operatorname{rank}_{\mathbb Q}R_A},
\]
and $A$ is rigid if and only if $R_A$ has full column rank.
\end{theorem}
\begin{proof}
The edge rows say that a kernel vector is constant on every connected component. The marked rows say that this constant is zero on every anchored component. Conversely, every assignment of a rational constant to each unanchored component, and zero to each anchored component, satisfies all rows. The indicator vectors of the unanchored components therefore form a basis of the kernel. In particular the nullity is $u_{\rm div}(A)$.

Intersecting with the binary cube permits exactly the independent choices zero or one on those components. The sign rule $1-2z_{\s}$ negates exactly the selected unanchored components. Theorem~\ref{ii:thm:subgroups} proves that each choice yields one binary subgroup-union mate, that every mate occurs, and that distinct choices have distinct inverses. This proves the bijection and the count.
\end{proof}

\begin{remark}[Scope and complexity]\label{v2:rem:rank-scope}
The theorem concerns the explicitly defined \emph{conductor-constraint} matrix. Its binary kernel is a cube because of the component-indicator basis; a general rational nullspace does not have this property. No compression to a lattice generated by the maximal subgroup generators is asserted here.

Given the prime factorization and an explicit binary valuation profile, the finite differences, cumulative sums, graph, and one selected inverse use $O(\omega(N)\tau(N))$ integer arithmetic operations. The rank can be read from the connected components without Gaussian elimination. This is an operation bound in the explicit valuation-grid representation, not a polynomial bit-time assertion for a succinct list of generators. Listing every mate necessarily includes the cost of writing all $2^{u_{\rm div}(A)}$ outputs.
\end{remark}
